% ==============================================================================
% METADATOS OFICIALES DEL PROYECTO DE INVESTIGACIÓN - ENTORNO SELLADO
% Cátedra: Proyecto de Investigación en Salud (EN 486) - FCS - UNSCH
% Investigadora Principal: Bach. Gresly Lucero Pariona Palomino
% Docente Asesor Titular: Dr. Manglio Aguirre Andrade
% ==============================================================================

% 1. Datos Institucionales
\newcommand{\metaUniversidad}{UNIVERSIDAD NACIONAL DE SAN CRISTÓBAL DE HUAMANGA}
\newcommand{\metaFacultad}{FACULTAD DE CIENCIAS DE LA SALUD}
\newcommand{\metaEscuela}{ESCUELA PROFESIONAL DE ENFERMERÍA}
\newcommand{\metaCatedra}{Proyecto de Investigación en Salud}
\newcommand{\metaCodigoCurso}{EN 486}

% 2. Título Oficial del Proyecto
\newcommand{\metaTitulo}{Calidad de Atención en Usuarios del Centro de Salud Belén, Ayacucho 2026}
\newcommand{\metaTituloMayus}{CALIDAD DE ATENCIÓN EN USUARIOS DEL CENTRO DE SALUD BELÉN, AYACUCHO 2026}
\newcommand{\metaTituloCorto}{Calidad de Atención - CS Belén 2026}
\newcommand{\metaTipoDocumento}{PROYECTO DE INVESTIGACIÓN CON ENFOQUE CUALITATIVO}
\newcommand{\metaGradoOptar}{LICENCIADA EN ENFERMERÍA}

% 3. Investigadora y Asesor
\newcommand{\metaAutora}{Gresly Lucero Pariona Palomino}
\newcommand{\metaAutoraMayus}{Gresly Lucero PARIONA PALOMINO}
\newcommand{\metaOrcidAutora}{0009-0008-5421-9872}
\newcommand{\metaAsesor}{Dr. Manglio Aguirre Andrade}
\newcommand{\metaAsesorMayus}{Dr. Manglio AGUIRRE ANDRADE}
\newcommand{\metaOrcidAsesor}{0000-0001-8234-567X}

% 4. Lugar y Año
\newcommand{\metaCiudad}{Ayacucho}
\newcommand{\metaPais}{Perú}
\newcommand{\metaLugarFecha}{AYACUCHO -- PERÚ}
\newcommand{\metaAno}{2026}

% 5. Alineación Normativa y Sanitaria Oficial
\newcommand{\metaLineaMinsa}{Línea 10: Sistemas y servicios de salud, acceso y cobertura universal (RM N° 424-2025/MINSA)}
\newcommand{\metaPrioridadDiresa}{Prioridad 10: Organización, gestión, calidad y accesibilidad de los servicios de salud (DIRESA Ayacucho 2025--2030)}
\newcommand{\metaEstablecimiento}{Centro de Salud Belén (Categoría I-3, Microred Huamanga, Red de Salud Huamanga)}

% 6. Metadatos Epistemológicos y Marco Teórico de Calidad
\newcommand{\metaEnfoque}{Cualitativo}
\newcommand{\metaDiseno}{Diseño Fenomenológico de Investigación (Edmund Husserl, Max van Manen, Roberto Hernández-Sampieri)}
\newcommand{\metaFenomenoCentral}{Calidad de Atención Asistencial}
\newcommand{\metaMarcoTeorico}{Proceso Interpersonal de Calidad (Avedis Donabedian) y Teoría del Cuidado Humano (Jean Watson)}
\newcommand{\metaPoblacion}{18,450 habitantes adscritos al Centro de Salud Belén}
\newcommand{\metaMuestra}{12 a 18 participantes informantes clave seleccionados por saturación teórica de categorías}
\newcommand{\metaRigor}{Criterios de rigor científico de Guba y Lincoln (1985)}
\newcommand{\metaSoftwareAnalisis}{ATLAS.ti v9}
